\documentclass[12pt,a4paper]{article}

% ====== ЯЗЫК И ШРИФТЫ (XeLaTeX) ======
\usepackage{fontspec}
\usepackage{polyglossia}
\setdefaultlanguage{russian}
\setotherlanguage{english}

\setmainfont{CMU Serif}
\setsansfont{CMU Sans Serif}
\setmonofont{CMU Typewriter Text}
\newfontfamily\cyrillicfont{CMU Serif}
\newfontfamily\cyrillicfontsf{CMU Sans Serif}
\newfontfamily\cyrillicfonttt{CMU Typewriter Text}

\usepackage{microtype}


% ====== МАТЕМАТИКА ======
\usepackage{amsmath, amssymb, amsthm, mathtools}
\newcommand{\nicefrac}[2]{#1/#2}

% Русские стандарты математической типографики
\renewcommand{\ge}{\geqslant}
\renewcommand{\le}{\leqslant}
\newcommand{\E}{\mathbb{E}}
\renewcommand{\P}{\mathbb{P}}
\newcommand{\F}{\mathcal{F}}

% ====== ГРАФИКА И TIKZ ======
\usepackage{tikz}
\usetikzlibrary{arrows.meta, positioning, automata, backgrounds, decorations.pathmorphing}

% ====== ОФОРМЛЕНИЕ СТРАНИЦЫ ======
\usepackage[left=2cm, right=2cm, top=2.5cm, bottom=2.5cm]{geometry}
\usepackage{parskip} % Красивые отступы между абзацами
\usepackage{xcolor}


% ====== ГИПЕРССЫЛКИ ======
\usepackage[colorlinks=true, linkcolor=blue!70!black, urlcolor=blue!70!black, citecolor=green!60!black]{hyperref}

% ====== TCOLORBOX (ОКРУЖЕНИЯ ДЛЯ ТЕОРЕМ И ЗАДАЧ) ======
\usepackage[most]{tcolorbox}
\tcbuselibrary{skins, breakable, theorems}

% Цветовая палитра "МЦНМО/ВШЭ"
\definecolor{defblue}{RGB}{20, 75, 140}
\definecolor{thmred}{RGB}{160, 30, 45}
\definecolor{probgreen}{RGB}{25, 110, 55}
\definecolor{warnpurple}{RGB}{130, 30, 130}
\definecolor{insightteal}{RGB}{0, 100, 100}
\definecolor{solgray}{RGB}{90, 90, 90}

% Определение
\newtcbtheorem[auto counter, number within=section]{definition}{Определение}{
    enhanced, breakable,
    colback=defblue!3, colframe=defblue,
    fonttitle=\bfseries\sffamily, title style={color=defblue},
    attach boxed title to top left={yshift=-2.5mm, xshift=5mm},
    boxed title style={sharp corners},
    boxrule=1pt, drop fuzzy shadow
}{def}

% Теорема
\newtcbtheorem[auto counter, number within=section]{theorem}{Теорема}{
    enhanced, breakable,
    colback=thmred!3, colframe=thmred,
    fonttitle=\bfseries\sffamily, title style={color=thmred},
    attach boxed title to top left={yshift=-2.5mm, xshift=5mm},
    boxed title style={sharp corners},
    boxrule=1pt, drop fuzzy shadow
}{thm}

% Задача
\newtcbtheorem[auto counter, number within=section]{problem}{Задача}{
    enhanced, breakable,
    colback=probgreen!3, colframe=probgreen,
    fonttitle=\bfseries\sffamily, title style={color=probgreen},
    attach boxed title to top left={yshift=-2.5mm, xshift=5mm},
    boxed title style={sharp corners},
    boxrule=1pt, drop fuzzy shadow
}{prob}

% Важное замечание / Парадокс
\newtcolorbox{paradox}[1][]{
    enhanced, breakable,
    colback=warnpurple!4, colframe=warnpurple!90!black,
    title={\textbf{Внимание: #1}},
    fonttitle=\bfseries\sffamily,
    coltitle=black,
    attach title to upper,
    after title={\par\smallskip},
    boxrule=1.5pt, arc=4pt
}

% Олимпиадный инсайт
\newtcolorbox{insight}[1][]{
    enhanced, breakable,
    colback=insightteal!3, colframe=insightteal, leftrule=4pt,
    boxrule=0pt, arc=0pt,
    title={\textbf{�� Олимпиадный инсайт: #1}},
    coltitle=insightteal, fonttitle=\bfseries\sffamily,
    attach title to upper,
    after title={\par\smallskip}
}

% Решение
\newenvironment{solution}{
    \begin{tcolorbox}[
        enhanced, breakable, blanker,
        borderline west={2.5pt}{0pt}{solgray},
        left=12pt, top=5pt, bottom=5pt,
        before skip=10pt, after skip=15pt
    ]
    \textbf{\textsf{РЕШЕНИЕ.}}\;
}{
    \hfill $\blacksquare$
    \end{tcolorbox}
}



\begin{document}

\begin{center}
    \vspace*{1.5cm}
    {\Huge \bfseries От Инвариантов к Мартингалам\par}
    \vspace{0.6cm}
    {\Large \itshape Искусство законов сохранения в вероятностных процессах\par}
    \vspace{0.8cm}
    \rule{0.6\textwidth}{1pt}\par
    \vspace{0.8cm}
    {\large \textbf{Калашников Александр}}\par
    \vspace{0.2cm}
    {\small Wildberries \& Russ, ФКН ВШЭ}
    \vspace{1.5cm}
\end{center}

\epigraph{\itshape «Случайность --- это лишь непознанная закономерность. Там, где дилетант видит хаос и бесконечные деревья исходов, настоящий математик находит мартингал».}{--- Джейсон Стейтем}

\tableofcontents
\newpage

\section{Введение: Эволюция закона сохранения}

В классической (детерминированной) олимпиадной математике фундаментом служат \textbf{инварианты} --- величины, которые строго не меняются после применения любых разрешённых правилами операций. 
Алгоритм решения прост: если начальное значение инварианта не равно конечному, цель недостижима (вспомните раскраски досок, остатки по модулю или алгебраическую сумму квадратов).

Но что происходит, когда в задачу вмешивается \textbf{случайность}? (Броски кубика, блуждания по графу, вероятностные турниры). Здесь строгих детерминированных инвариантов нет: любая величина на следующем шаге может вырасти, а может и упасть. Однако мы можем потребовать, чтобы величина сохранялась \textit{в среднем}. 

\begin{center}
    \large
    \textbf{Инвариант} $\implies$ сохраняется \textbf{абсолютно}.\\
    \textbf{Мартингал} $\implies$ сохраняется \textbf{в математическом ожидании}.
\end{center}

Это концептуальный скачок, который позволяет решать тяжелейшие вероятностные задачи без составления бесконечных систем линейных алгебраических уравнений (СЛАУ) или суммирования рядов.

\section{Первое касание: Инварианты матожидания}

Давайте посмотрим на задачу, где результат каждого шага случаен, но математическое ожидание системы ведёт себя жёстко и предсказуемо.

\begin{problem}{Игра Незнайки (ВсОШ)}{neznayka}
Незнайка $2026$ раз подкидывает волшебную монетку. Изначально вероятность выпадения орла $p_1=0.5$. 
\begin{itemize}
    \item Если выпадает \textbf{орёл}, вероятность орла в будущем падает на $0.1$, и Незнайка выигрывает $1$ тугрик. 
    \item Если выпадает \textbf{решка}, вероятность орла растёт на $0.1$, а выигрыш равен $0$. 
\end{itemize}
Сколько тугриков в среднем он выиграет за всю игру? (Гарантируется, что $p_k \in [0, 1]$).
\end{problem}

\begin{solution}
Обозначим состояние монетки перед $k$-м броском через $X_k = 10 \cdot p_k$. Изначально $X_1 = 5$.
Посчитаем \textbf{условное} математическое ожидание состояния на следующем шаге, зная, что сейчас $X_k = x$:
\begin{align*}
\E[X_{k+1} \mid X_k = x] &= \underbrace{\left(\frac{x}{10}\right)}_{\P(\text{орёл})} \cdot (x-1) \;+\; \underbrace{\left(\frac{10-x}{10}\right)}_{\P(\text{решка})} \cdot (x+1) \\
&= \frac{x^2 - x + 10x + 10 - x^2 - x}{10} = \frac{8x + 10}{10} = 0.8x + 1.
\end{align*}
(Заметим, что формула работает идеально даже на границах: при $x=0$ матожидание равно $1$, при $x=10$ равно $9$).

Используя свойство линейности, перейдём к безусловному матожиданию: $m_k = \E[X_k]$. 
Получаем рекурренту: $m_{k+1} = 0.8m_k + 1$. \\
Так как $m_1 = 5$, проверяем следующий шаг: $m_2 = 0.8 \cdot 5 + 1 = 5$. По индукции $m_k = 5$ для любого $k$.

Значит, \textbf{ожидаемая} вероятность выпадения орла на любом шаге неизменна и равна $0.5$. Так как выигрыш равен 1 только при выпадении орла, ожидаемый куш за $2026$ бросков составит $2026 \cdot 0.5 = \mathbf{1013}$ тугриков.
\end{solution}

\section{Строгая теория: Мартингалы и Теорема Дуба}

В задаче про Незнайку сохранялось \textit{безусловное} ожидание. Но истинная мощь мартингалов раскрывается тогда, когда сохраняется \textbf{условное} ожидание — то есть когда локально, на каждом шаге, игра абсолютно честна.

\begin{definition}{Мартингал}{mart}
Пусть $\{M_n\}_{n \ge 0}$ --- последовательность случайных величин. Через $\F_n$ обозначается \textit{фильтрация} --- вся накопленная информация о процессе до шага $n$ включительно. Процесс называется \textbf{мартингалом}, если:
\begin{enumerate}
    \item $\E[|M_n|] < \infty$ (математическое ожидание существует и конечно).
    \item Значение $M_n$ полностью определяется известной историей $\F_n$.
    \item \textbf{Свойство честной игры:} $\E[M_{n+1} \mid \F_n] = M_n$.
\end{enumerate}
\end{definition}

Смысл третьего пункта: \textit{«Даже зная всю историю процесса до текущей секунды, моё лучшее математическое предсказание капитала на завтра — это мой капитал сегодня»}.

Из определения по индукции следует, что $\E[M_n] = \E[M_0]$ для любого \textbf{фиксированного} момента $n$. Но в олимпиадных задачах игра почти всегда останавливается в \textbf{случайный} момент $\tau$ (достижение цели, банкротство, появление нужной строки). Сохранится ли инвариант?

\begin{paradox}[Ловушка бесконечного капитала (Стратегия Мартингейл)]
Самая известная система игры в рулетку: ставь 1 рубль на красное; если проиграл — ставь 2; проиграл — ставь 4, и так до первой победы. Как только выигрываешь, забираешь чистую прибыль (1 рубль) и уходишь (это и есть момент остановки $\tau$).

Каждая ставка честная (ожидание $0$), значит ваш капитал $M_n$ — мартингал.
Очевидно, игра конечна: вероятность вечно проигрывать равна нулю ($\P(\tau < \infty) = 1$).
В момент остановки $\tau$ ваш баланс \textbf{строго равен} $+1$ рублю. То есть $M_\tau \equiv 1$, значит $\E[M_\tau] = 1$. Но изначально у вас было 0 рублей ($\E[M_0] = 0$). Мы математически доказали, что $\mathbf{1 = 0}$! Где ошибка?

\textbf{Разгадка:} В этой стратегии возможный минус на балансе до выигрыша ничем не ограничен (игрок может уйти в бесконечные долги). Бездушная математика запрещает приравнивать матожидания для процессов, которые могут неограниченно падать вниз. Нам нужна теорема, которая защитит нас от «генерации денег из воздуха».
\end{paradox}

\begin{theorem}{Теорема Дуба об опциональной остановке}{ost}
Пусть $\{M_n\}$ --- мартингал, а $\tau$ --- марковский момент остановки (решение об остановке не использует взгляд в будущее). Равенство $\mathbf{\E[M_\tau] = \E[M_0]}$ гарантированно выполняется, если выполнено \textbf{хотя бы одно} из условий регулярности:
\begin{enumerate}
    \item Время игры строго ограничено ($\tau \le C$).
    \item Ожидаемое время конечно ($\E[\tau] < \infty$), а «прыжки» ограничены ($|M_{n+1} - M_n| \le K$).
    \item \textbf{Практическое:} Значения мартингала ограничены вплоть до остановки ($|M_{n \wedge \tau}| \le C$).
\end{enumerate}
\end{theorem}

\begin{insight}[Как оформлять решение на чистовик]
На школьных олимпиадах строгая проверка условий Теоремы Дуба почти никогда не требует сложных выкладок. Если процесс происходит на конечном графе или капитал ограничен (например, колеблется от $0$ до $N$), смело пишите: 
\textit{«Так как состояния системы ограничены, и из любого состояния вероятность достичь конца за $N$ шагов $\ge \varepsilon > 0$, игра завершится с вероятностью 1. Условия Теоремы об остановке выполнены.»} Жюри поставит полный балл.
\end{insight}

\section{Пять видов мартингалов (Тяжелая артиллерия)}

Существует 5 классов мартингалов, которые покрывают 99\% вероятностных задач высшего уровня. Поймёте их — и СЛАУ вам больше не понадобятся.

\subsection{Линейные мартингалы (Симметричные игры)}

\begin{problem}{Турнир гладиаторов (ВсОШ по ИИ)}{gladiator}
Гладиаторы армии Тимофея (сумма сил $S_T = 55$) и армии Максима ($S_M = 30$) сошлись в турнире. В каждом бою случайно выбираются два любых бойца (с силами $x$ и $y$). Боец с силой $x$ побеждает $y$ с вероятностью $\frac{x}{x+y}$ и забирает его силу себе. Проигравший умирает. Найти вероятность победы армии Максима.
\end{problem}

\begin{solution}
Рассмотрим суммарную силу армии Тимофея $S_T$. Найдём ожидаемое изменение её силы $\Delta S_T$ за один бой между $x \in T$ и $y \in M$:
\[ \E[\Delta S_T \mid \F_n] = \underbrace{(+y)}_{\text{Тимофей забрал силу } y} \cdot \left(\frac{x}{x+y}\right) \;\; + \;\; \underbrace{(-x)}_{\text{Тимофей потерял бойца } x} \cdot \left(\frac{y}{x+y}\right) = \frac{xy - xy}{x+y} = 0. \]
\textbf{Поразительный факт!} Неважно, по какому алгоритму генералы выбирают бойцов на арену. Для любых $x$ и $y$ матожидание изменения равно нулю. Процесс $S_T$ --- \textbf{строгий мартингал}.

Игра конечна (в каждом бою умирает 1 человек), состояния ограничены от 0 до 85. В момент окончания $\tau$ сила Тимофея равна 85 (вероятность $P_T$) или 0 (вероятность $1-P_T$). По Теореме Дуба:
\[ \E[S_\tau] = S_0 \implies 85 \cdot P_T + 0 \cdot (1-P_T) = 55. \]
Отсюда $P_T = \nicefrac{55}{85} = \nicefrac{11}{17}$. Вероятность победы Максима равна $\mathbf{\nicefrac{6}{17}}$.
\end{solution}

\subsection{Экспоненциальные мартингалы (Асимметрия и Искривление)}

Что делать, если игра \textbf{нечестная}? Например, игрок делает шаг $+1$ с вероятностью $p$ и $-1$ с вероятностью $q$ ($p \neq q$). В этом случае капитал $X_n$ имеет тренд (он убывает или растет), и линейный мартингал ломается. Школьный метод требует составления жуткой рекурренты 2-го порядка. Олимпиадный метод — \textit{искривить} метрику пространства.

\begin{center}
\begin{tikzpicture}[>=Stealth, node distance=2.2cm, 
    state/.style={circle, draw=defblue, thick, fill=defblue!5, minimum size=1.2cm, font=\bfseries}]
    
    \node[state] (s0) {$0$};
    \node[state] (s1) [right=of s0] {$1$};
    \node[state] (s2) [right=of s1] {$2$};
    \node (dots) [right=of s2] {\Large $\dots$};
    \node[state] (sN) [right=of dots] {$N$};

    \path[->, thick, color=probgreen] 
        (s1) edge[bend left=35] node[above] {$p$} (s2)
        (s2) edge[bend left=35] node[above] {$p$} (dots)
        (dots) edge[bend left=35] node[above] {$p$} (sN);
        
    \path[->, thick, color=thmred] 
        (sN) edge[bend left=35] node[below] {$q$} (dots)
        (dots) edge[bend left=35] node[below] {$q$} (s2)
        (s2) edge[bend left=35] node[below] {$q$} (s1)
        (s1) edge[bend left=35] node[below] {$q$} (s0);
\end{tikzpicture}
\end{center}

\begin{problem}{Несправедливое разорение}{}
У игрока $A=10$ руб, он хочет достичь $N=100$ руб. На каждом шаге он выигрывает 1 руб с вероятностью $p=0.48$ и теряет с $q=0.52$. Найти вероятность $P_{win}$ того, что он доберётся до цели.
\end{problem}

\begin{solution}
Магическим инвариантом для несимметричного блуждания является функция $M_n = \left(\frac{q}{p}\right)^{X_n}$. 
Проверим её условное матожидание:
\[ \E[M_{n+1} \mid X_n] = p \cdot \left(\frac{q}{p}\right)^{X_n+1} + q \cdot \left(\frac{q}{p}\right)^{X_n-1} = \left(\frac{q}{p}\right)^{X_n} \cdot \left[ p \frac{q}{p} + q \frac{p}{q} \right] = \left(\frac{q}{p}\right)^{X_n} (q + p) = M_n. \]
Экспонента идеально скомпенсировала дрейф! Применяем Теорему Дуба: $\E[M_\tau] = M_0$.
\[ P_{win} \cdot \left(\frac{q}{p}\right)^N + (1 - P_{win}) \cdot \left(\frac{q}{p}\right)^0 = \left(\frac{q}{p}\right)^A. \]
Так как любое число в нулевой степени --- это 1, получаем готовую формулу:
\[ P_{win} = \mathbf{\frac{(q/p)^A - 1}{(q/p)^N - 1}}. \]
Подставляем: $q/p = \nicefrac{0.52}{0.48} \approx 1.083$. Вероятность победы ничтожна: $P_{win} \approx \frac{(1.083)^{10} - 1}{(1.083)^{100} - 1} \approx 0.0004$.
\end{solution}

\subsection{Квадратичные мартингалы (Поиск времени)}

Линейные и экспоненциальные мартингалы находят \textit{вероятности} исходов, но "схлопывают" время. Чтобы найти ожидаемое время $\E[\tau]$, нужен мартингал, который явно вычитает $n$.

\begin{problem}{Время разорения}{}
Два игрока с капиталами $A$ и $B$ играют в честную орлянку со ставкой 1. Сколько бросков монеты в среднем продлится игра до разорения одного из них?
\end{problem}

\begin{solution}
Пусть $X_n$ --- капитал первого игрока. Мы знаем, что $X_n$ — линейный мартингал. А как ведёт себя его квадрат?
\[ \E[X_{n+1}^2 \mid \F_n] = \frac{1}{2}(X_n + 1)^2 + \frac{1}{2}(X_n - 1)^2 = X_n^2 + 1. \]
Квадрат капитала в среднем растёт строго на $1$ каждую секунду. 
Следовательно, функция $M_n = X_n^2 - n$ обязана быть мартингалом! Проверим:
\[ \E[M_{n+1} \mid \F_n] = \E[X_{n+1}^2 - (n+1) \mid \F_n] = (X_n^2 + 1) - n - 1 = X_n^2 - n = M_n. \]
Применяем Теорему Дуба: $\E[M_\tau] = M_0 \implies \E[X_\tau^2 - \tau] = A^2 \implies \E[\tau] = \E[X_\tau^2] - A^2$.

В конце игры капитал $X_\tau$ равен либо $A+B$ (с вероятностью $\frac{A}{A+B}$), либо $0$.
\[ \E[X_\tau^2] = (A+B)^2 \cdot \frac{A}{A+B} + 0^2 \cdot \frac{B}{A+B} = A(A+B). \]
Подставляем в формулу: $\E[\tau] = A(A+B) - A^2 = \mathbf{AB}$. Время честной игры в точности равно произведению стартовых капиталов!
\end{solution}

\subsection{Сдвиговые мартингалы и Тождество Вальда}

А что если мы суммируем не фиксированное число случайных величин, а \textbf{случайное} число? И при этом игра нечестная (ожидаемое смещение за 1 шаг $\mu \neq 0$)? Здесь работает сдвиговый мартингал $M_n = X_n - \mu \cdot n$.

\begin{theorem}{Тождество Вальда (Wald's Equation)}{wald}
Пусть $Y_1, Y_2, \ldots$ --- независимые одинаково распределённые величины с матожиданием $\E[Y]$. Пусть $N$ --- случайный момент остановки, причём $\E[N] < \infty$. Тогда математическое ожидание случайной суммы $S_N = Y_1 + Y_2 + \ldots + Y_N$ равно:
\[ \mathbf{\E[S_N] = \E[N] \cdot \E[Y]} \]
\end{theorem}

\begin{proof}
Рассмотрим процесс $M_n = S_n - n \E[Y]$. Проверим мартингал:
\[ \E[M_{n+1} \mid \F_n] = \E[S_n + Y_{n+1} - (n+1)\E[Y] \mid \F_n] = S_n + \E[Y] - n\E[Y] - \E[Y] = M_n. \]
По Теореме Дуба $\E[M_N] = M_0 = 0 \implies \E[S_N - N \E[Y]] = 0 \implies \E[S_N] = \E[N]\E[Y]$.
\end{proof}

\begin{problem}{Броски кубика (Парадокс остановки)}{}
Мы бросаем стандартный кубик до тех пор, пока не выпадет шестёрка (включительно). Какова ожидаемая \textbf{сумма} всех выпавших очков за эту серию бросков?
\end{problem}
\begin{solution}
Количество бросков $N \sim \text{Geom}(1/6)$, значит $\E[N] = 6$. Матожидание одного броска (любого) равно $\E[Y] = 3.5$. По тождеству Вальда, матожидание суммы: $\E[S_N] = 6 \cdot 3.5 = \mathbf{21}$. 
\end{solution}

\begin{paradox}[Вопрос от отличника]
«Но ведь последний бросок в этой серии \textbf{всегда} равен 6! Значит, его ожидание не 3.5, а 6. Разве мы можем просто умножить $\E[N]$ на безусловное $\E[Y] = 3.5$?» 

\textbf{Разгадка:} Давайте посчитаем вручную. Сумма равна $(N-1)$ бросков без шестерок плюс последняя шестерка. $\E[N-1] = 5$. Ожидание броска при условии, что это не шестерка: $(1+2+3+4+5)/5 = 3$. 
Итого: $5 \cdot 3 + 6 = \mathbf{21}$. Математика Вальда абсолютно безупречна! Мартингал элегантно прячет зависимость последнего элемента внутрь случайности самой длины суммы.
\end{paradox}

\subsection{Дробные мартингалы (Процессы с памятью)}

Иногда инварианты работают даже там, где вероятности зависят от предыстории (памяти).

\begin{problem}{Урна Пойа (Pólya Urn)}{polya}
В урне $R_0$ красных и $B_0$ синих шаров. На каждом шаге мы наугад достаём шар, и возвращаем его обратно, добавляя ещё $c$ шаров \textbf{такого же цвета}. Какова ожидаемая доля красных шаров после $N$ шагов?
\end{problem}

\begin{solution}
Пусть $S_n = R_0 + B_0 + n \cdot c$ --- детерминированное число шаров на шаге $n$. Пусть $R_n$ --- количество красных. Найдём математическое ожидание доли $M_n = \nicefrac{R_n}{S_n}$ на следующем шаге:
\[
\E[R_{n+1} \mid \F_n] = \underbrace{M_n \cdot (R_n + c)}_{\text{добавили красные}} + \underbrace{(1 - M_n) \cdot R_n}_{\text{добавили синие}} = R_n + c \cdot M_n.
\]
Переходим к доле:
\[
\E[M_{n+1} \mid \F_n] = \frac{R_n + c \left(\frac{R_n}{S_n}\right)}{S_n + c} = \frac{R_n \left(\frac{S_n + c}{S_n}\right)}{S_n + c} = \frac{R_n}{S_n} = M_n.
\]
Доля $M_n$ является строгим мартингалом! Следовательно, ожидаемая доля красных шаров на \textbf{любом} шаге, даже спустя миллион итераций, равна стартовой пропорции $\mathbf{\frac{R_0}{R_0 + B_0}}$. (Более того, по теореме сходимости, при $n \to \infty$ эта доля сходится к случайной величине, имеющей Бета-распределение).
\end{solution}

\section{Жемчужина: Метод Мартина-Лёфа (Трюк с Казино)}

В олимпиадном программировании (задачи на строки/NLP/автоматы) составить СЛАУ для генерации длинной строки невозможно. На помощь приходит гениальный трюк (алгоритм Conway / ABR).

\begin{problem}{Обезьяна и АБРАКАДАБРА}{}
Обезьяна равновероятно нажимает клавиши (33 буквы алфавита). Найти матожидание количества нажатий, через которое впервые напечатается слово \textbf{АБРАКАДАБРА} (длина 11).
\end{problem}

\begin{solution}
Представим \textbf{идеально честное казино}, в которое \textit{перед каждым} нажатием клавиши заходит новый Игрок. У Игрока есть ровно 1 рубль.
\begin{itemize}
    \item Игрок ставит 1 рубль на первую букву «\textbf{А}». Не угадал — ушёл. Угадал (шанс 1/33) — казино выплачивает ему $33$ рубля.
    \item Свой выигрыш Игрок ставит на вторую букву «\textbf{Б}». Угадал — получает $33^2$.
    \item Так он продолжает, пока не пройдет всё слово, сорвав джекпот $33^{11}$.
\end{itemize}

Поскольку игра честная, баланс казино $M_n$ (сборы минус выплаты) является \textbf{мартингалом}.
Пусть $\tau$ — момент, когда слово напечатано. По Теореме Дуба: $\E[M_\tau] = 0 \implies \E[\text{Сборы}] = \E[\text{Выплаты}]$.

В каждый из $\tau$ тактов заходил 1 игрок с 1 рублем. Казино собрало $\E[\tau]$ рублей!
А кому казино выплачивает деньги в момент $\tau$? Смотрим на конец строки:
\dots \textbf{А Б Р А К А Д А Б Р А}.
\begin{enumerate}
    \item Игрок, зашедший 11 ходов назад, угадал всё слово! Ему платим \textbf{$33^{11}$}.
    \item Игрок, зашедший 4 хода назад, ставил на «А», «Б», «Р», «А». И они совпали! (Суффикс длины 4 совпадает с префиксом длины 4). Ему платим \textbf{$33^4$}.
    \item Игрок, зашедший на последнем такте, ставил на «А». И последняя буква — «А». Ему платим \textbf{$33^1$}.
\end{enumerate}

Суммируя долги, получаем точный ответ без решения матрицы $11 \times 11$:
\[ \mathbf{\E[\tau] = 33^{11} + 33^4 + 33^1}. \]

\textbf{Алгоритм «Казино»:} Время генерации слова равно $\sum K^m$, где $K$ --- алфавит, а $m$ --- длины префиксов слова, которые одновременно являются его суффиксами.
\end{solution}

\begin{insight}[Парадокс Конвея (Игры Пенни)]
Чего мы дождемся быстрее при бросках монеты: комбинации \textbf{ОРО} (Орел-Решка-Орел) или \textbf{РОО} (Решка-Орел-Орел)? Применим алгоритм Казино (алфавит $K=2$).

\textbf{Для ОРО:} Слово совпадает само с собой по длине 3 (ОРО) и по длине 1 (О). 
$\E[\tau_{ОРО}] = 2^3 + 2^1 = \mathbf{10}$ бросков.

\textbf{Для РОО:} Совпадает с собой только по длине 3 (префикс Р $\neq$ суффикс О). 
$\E[\tau_{РОО}] = 2^3 = \mathbf{8}$ бросков.

Абсолютно поразительно! Комбинация \textbf{РОО} выпадает в среднем быстрее, чем \textbf{ОРО}, хотя их вероятности совершенно одинаковы. 
\end{insight}

\newpage

\section{Шпаргалка Олимпиадника (Cheat Sheet)}

Перед финалом ВсОШ или перечневой олимпиадой прогоняйте любую вероятностную задачу по этому сканеру:

\begin{tcolorbox}[colback=yellow!5, colframe=orange!80!black, title=\textbf{Как выбрать правильный мартингал?}, fonttitle=\bfseries]
\begin{enumerate}
    \item \textbf{Найти вероятность победы (Симметричная игра $p=q$):}\\
    Ищем классический линейный мартингал $X_n$ (сумма ресурсов). Решаем: $X_{start} = P \cdot X_{win} + (1-P) \cdot X_{loss}$.
    
    \item \textbf{Найти вероятность победы (Асимметричная игра $p \neq q$):}\\
    Используем Экспоненциальный мартингал Муавра $(q/p)^{X_n}$. Он мгновенно линеаризует процесс.
    
    \item \textbf{Найти ожидаемое время игры (Симметричная игра):}\\
    Если $X_n$ меняется на $\pm 1$, то $X_n^2 - n$ является мартингалом. Приравниваем $X_0^2 = \E[X_\tau^2] - \E[\tau]$. 
    
    \item \textbf{Найти ожидаемое время (С дрейфом $\mu$ / Тождество Вальда):}\\
    Если матожидание шага $\mu \neq 0$, используйте мартингал $X_n - \mu n$. 
    Отсюда $\E[\tau] = (\E[X_\tau] - X_0) / \mu$.
    
    \item \textbf{Найти время генерации слова (Строки, NLP, Графы):}\\
    Включаем логику «Казино». Ищем все перекрытия слова (префиксы, равные суффиксам). Ответ всегда имеет вид суммы степеней алфавита.
\end{enumerate}
\end{tcolorbox}

\vspace{0.8cm}
\begin{center}
\textit{Удачи! Помните: умение находить строгий порядок внутри кажущегося случайным хаоса --- абсолютный признак высшего математического мастерства.}
\end{center}

\end{document}